% ---------------------------------------------------------------------------
% The birthmark. (a) The transverse plane of our event to scale: the primary
% vertex, the invisible walk of the Bbar0 to its decay vertex 9.9 mm away,
% the eight charged tracks that leave it, the D0 walking another 5.1 mm, and
% its three tracks. (b) Not to scale: the impact parameter of one track, the
% distance by which its line misses the primary vertex, and why it does not
% grow with the boost. (c) The numbers of the ensemble.
% ---------------------------------------------------------------------------
\begin{tikzpicture}[x=1cm,y=1cm,
  every node/.style={font=\footnotesize},
  lab/.style={font=\scriptsize,align=center},
  head/.style={font=\small\bfseries},
  trk/.style={red!75!black,thick},
  ctrk/.style={orange!80!black,thick},
  walk/.style={red!80!black,line width=1.4pt,dashed},
  cwalk/.style={orange!80!black,line width=1.2pt,dashed}]
% ---------------- (a) to scale: 0.32 cm per mm --------------------------------------------------
\begin{scope}
\node[head,anchor=west] at (-6.2,2.3) {(a) our event in the transverse plane ($0.32\,$cm per mm)};
\draw[-{Stealth}] (-6.0,-2.6) -- (1.0,-2.6) node[below left,lab] {$x$ [mm]};
\foreach \x/\t in {-5.44/-17,-3.84/-12,-2.24/-7,-0.64/-2,0.96/3} {\draw (\x,-2.66) -- (\x,-2.54);}
\foreach \x/\t in {-4.8/-15,-3.2/-10,-1.6/-5,0/0} {\draw (\x,-2.7) -- (\x,-2.5); \node[lab,anchor=north] at (\x,-2.72) {\t};}
\draw[-{Stealth}] (0.9,-2.6) -- (0.9,1.6) node[left,lab] {$y$ [mm]};
\foreach \y/\t in {-1.92/-6,-0.64/-2,0.64/2} {\draw (0.84,\y) -- (0.96,\y); \node[lab,anchor=west] at (0.98,\y) {\t};}
\node[star,star points=5,fill=magenta,inner sep=1.6pt] (pv) at (0,0) {};
\node[lab,anchor=south west] at (0.05,0.08) {primary vertex};
\coordinate (sv) at (-2.92,-1.26);   % (-9.13, -3.93) mm
\coordinate (tv) at (-4.44,-1.87);   % (-13.86, -5.86) mm
\draw[walk] (pv) -- (sv); \node[lab,text=red!80!black,anchor=south east] at (-1.2,-0.45) {$\bar B^0$: no track,\\ $L_{xy}=9.9\,\mathrm{mm}$};
\fill[red!80!black] (sv) circle (1.1mm); \node[lab,text=red!80!black,anchor=north] at (-2.92,-1.42) {decay vertex\\ $(-9.1,-3.9)$ mm};
\foreach \a in {-176,-168,-161,-152,-143,-128,-110,-100} {\draw[trk] (sv) -- ++(\a:3.2);}
\draw[cwalk] (sv) -- (tv); \fill[orange!80!black] (tv) circle (1.0mm);
\node[lab,text=orange!80!black,anchor=north] at (-4.44,-2.0) {$D^0$: $+5.1\,\mathrm{mm}$};
\foreach \a in {-170,-158,-150} {\draw[ctrk] (tv) -- ++(\a:1.7);}
\node[lab,anchor=north west,text width=6.9cm,align=left,text=black!65] at (-6.2,-3.1) {eight charged tracks from the $\bar B^0$ and its charm, all of them missing the primary vertex by more than $100\,\mu$m, the largest by $1.64\,\mathrm{mm}$; nothing of the walk itself is seen};
\end{scope}
% ---------------- (b) the impact parameter, not to scale ------------------------------------------
\begin{scope}[xshift=9.2cm]
\node[head,anchor=west] at (-0.6,2.3) {(b) the impact parameter $d_0$};
\node[star,star points=5,fill=magenta,inner sep=1.6pt] (pv2) at (0,0) {};
\node[lab,anchor=north] at (0,-0.2) {primary vertex};
\coordinate (sv2) at (3.6,1.1);
\draw[walk] (pv2) -- (sv2); \node[lab,text=red!80!black,anchor=south] at (1.8,0.7) {$\vec L$ of the B};
\fill[red!80!black] (sv2) circle (1.1mm);
\draw[trk] (sv2) -- ++(10:2.4) node[right,lab,text=red!75!black] {a track, $\vec p_{\mathrm{T}}$};
\draw[trk,dotted] (sv2) -- ++(190:4.2);
\coordinate (end2) at ($(sv2)+(190:4.2)$); \coordinate (foot) at ($(sv2)!(pv2)!(end2)$);
\draw[blue!70!black,thick,{Stealth}-{Stealth}] (pv2) -- (foot) node[midway,right,lab,text=blue!70!black] {$d_0$};
\draw[black!60] (sv2) ++(190:0.8) arc (190:197:0.8); \node[lab] at (2.5,1.35) {$\theta$};
\node[lab,anchor=north west,text width=6.2cm,align=left] at (-0.6,-0.5) {$d_0=\dfrac{|x_v\,p_y-y_v\,p_x|}{\pt}=L_{xy}\sin\theta$, with $\theta$ the angle between the track and the B direction. The boost makes $L$ larger and $\theta$ smaller by the same factor: $d_0$ stays of the order of $c\tau_0\approx0.5\,\mathrm{mm}$ at any jet energy. That is why the tracker, not the calorimeter, finds the B};
\end{scope}
% ---------------- (c) the numbers ------------------------------------------------------------------
\begin{scope}[xshift=16.6cm]
\node[head,anchor=west] at (-0.4,2.3) {(c) 400 decays of the same hadrons};
\node[lab,anchor=north west,text width=5.4cm,align=left] at (-0.4,1.9) {%
\textbf{the $b$'s $\bar B^0$:} $4.85$ charged tracks per decay, $3.0$ of them with $|d_0|>100\,\mu$m; the median $|d_0|$ is $166\,\mu$m, the mean $362$, the largest track of a decay $874\,\mu$m on average\\[3pt]
\textbf{\PYTHIA on the same hadrons:} $4.96$ tracks, $3.0$ above $100\,\mu$m, mean $|d_0|$ $330\,\mu$m\\[3pt]
\textbf{our event:} 8 tracks, 8 above $100\,\mu$m\\[3pt]
\textbf{the \CMSSW $B^-$:} $L=7.83$, $L_{xy}=3.13\,\mathrm{mm}$\\[3pt]
the charm hadron walks on: $D^0$ $c\tau_0=0.123\,\mathrm{mm}$, in our event $11.0\,\mathrm{mm}$ at $\gamma\beta=31$; a third vertex for the $b$ tag of the later chapters};
\end{scope}
\end{tikzpicture}
